\documentclass[11pt,a4paper]{article}
\usepackage{fontspec,amsmath,amssymb,amsthm,microtype}
\setmainfont{lmroman10-regular.otf}[BoldFont=lmroman10-bold.otf,ItalicFont=lmroman10-italic.otf,BoldItalicFont=lmroman10-bolditalic.otf]
\setmonofont{lmmono10-regular.otf}[Scale=MatchLowercase]
\usepackage[margin=25mm]{geometry}
\usepackage[unicode,hidelinks]{hyperref}
\setlength{\emergencystretch}{3em}
\setlength{\parskip}{0.4em}
\setlength{\parindent}{0pt}
\allowdisplaybreaks[2]
\raggedbottom
\title{\texorpdfstring{One actual KLMC2 output law for every initial coupling}{One actual KLMC2 output law for every initial coupling}}
\author{LangevinLean mathematical exposition}
\date{}
\begin{document}\maketitle
\noindent Source module: \texttt{Kinetic\allowbreak{}Langevin/\allowbreak{}KLMC2\allowbreak{}Product\allowbreak{}Law.\allowbreak{}lean}. The accompanying HTML pairs this mathematical exposition with the exact Lean statements and records the scope of the formalization.


The stochastic analysis may choose any initial transport coupling. The implemented sampler has one fixed law, defined independently of this choice. This module proves the two constructions coincide exactly, including the random Hessian and all four source noises. It supplies the law identity needed to pass a uniform coupling estimate to the true Wasserstein infimum.


\section{\texorpdfstring{Lift the actual four-\allowbreak{}noise source}{Lift the actual four-noise source}}


The initial coupling and its independent source extension are probability measures. Each square-integrable original noise coordinate is composed with the source projection and represented in L\ensuremath{{}^{2}}. Intersecting the four representative identities proves that these lifted variables retain precisely the original four-coordinate noise value.



{\raggedright\noindent Lean declarations: \texttt{second\allowbreak{}Order\allowbreak{}Coupling\allowbreak{}Probability}, \texttt{second\allowbreak{}Order\allowbreak{}Extension\allowbreak{}Probability}, \texttt{extension\allowbreak{}KLMC2\allowbreak{}Noise}, \texttt{extension\allowbreak{}KLMC2\allowbreak{}Noise\_\allowbreak{}coe}.\par}

\section{\texorpdfstring{Define the fixed-\allowbreak{}source actual sampler law}{Define the fixed-source actual sampler law}}


The pointwise original KLMC2 recursion starts at the common source velocity and the independent initial position. Its terminal position defines the output map and its product-source pushforward defines the actual law. The force and Hessian are continuous, so the complete four-noise step is measurable. Induction over the original recursion proves the output map is almost everywhere measurable from the square-integrable source data.



{\raggedright\noindent Lean declarations: \texttt{product\allowbreak{}KLMC2\allowbreak{}Map}, \texttt{product\allowbreak{}KLMC2\allowbreak{}Law}, \texttt{product\allowbreak{}KLMC2\allowbreak{}Map\_\allowbreak{}aemeasurable}.\par}

\section{\texorpdfstring{Identify the L\ensuremath{{}^{2}} analysis with that single law}{Identify the L² analysis with that single law}}


The pointwise correspondence theorem for the genuine random-Hessian L\ensuremath{{}^{2}} chain identifies every iterate with the original vector recursion. The initial velocity, initial coupling coordinate and all countably many four-noise values have exactly their source representatives. Substituting these identities proves the L\ensuremath{{}^{2}} output equals the original fixed-source output map composed with the initial/\allowbreak{}source projection. That projection has exactly the independent initial-position/\allowbreak{}source product law. Pushforward composition therefore proves the actual output law is independent of the chosen transport coupling. Its measurability also confirms it is a probability measure.



{\raggedright\noindent Lean declarations: \texttt{product\_\allowbreak{}klmc2\allowbreak{}Chain\_\allowbreak{}coe}, \texttt{product\_\allowbreak{}klmc2\allowbreak{}Chain\_\allowbreak{}output\allowbreak{}Law}, \texttt{product\allowbreak{}KLMC2\allowbreak{}Law\_\allowbreak{}probability}.\par}

\end{document}
